\documentclass[12pt]{article}
\usepackage[utf8]{inputenc}
\usepackage[T1]{fontenc}
\usepackage{hyperref}
\title{Yamphu Language Resources for K-12 Curriculum Development}
\author{Compiled by Claude Sonnet 5.5 (claude-sonnet-5-5) for Andr\'as Kornai}
\date{}
\begin{document}
\maketitle

\section{Introduction}
Yamphu (ISO 639-3 ybi; Glottocode yamp1242) is a Kiranti language of the Sino-Tibetan family, spoken in Nepal by the Yamphu Rai. The 2011 census gives 9,200 native speakers. The Endangered Languages Catalogue lists it as threatened. The searches found one major scholarly description and no school materials.

\section{Primary Educational Resources}
No primer, textbook or school programme for Yamphu turned up in the searches.

\section{Linguistic and Reference Materials}
\begin{itemize}
\item Roland Rutgers, \emph{Yamphu: Grammar, Texts \&amp; Lexicon} (Leiden: CNWS, 1998, 632 pp., Languages of the Greater Himalayan Region 2). It covers the Hedangna dialect and includes oral texts and a lexicon. Glottolog reference: \url{https://glottolog.org/resource/reference/id/12930}
\item Bookseller listing of the same volume: \url{https://www.boekwinkeltjes.nl/b/230893072/}
\item STEDT source entry: \url{https://stedt.berkeley.edu/~stedt-cgi/rootcanal.pl/source/RR-Yamphu}
\item Wikipedia article on the language: \url{https://en.wikipedia.org/wiki/Yamphu_language}
\item Endangered Languages bibliography: \url{https://endangeredlanguages.com/lang/3216/bibliography}
\end{itemize}

\section{Religious and Community Resources}
No religious or community-organisation resources turned up in the searches.

\section{Digital and Audio Resources}
No audio, video or digital resources turned up in the searches. The Rutgers volume is the only lexicon found, and no separate dictionary was found.

\section{Curriculum Development Context}
Yamphu has one detailed description (1998) and little else. A curriculum would need new community-led materials and orthography decisions, starting from Rutgers's texts and lexicon.

\end{document}